%% 2-1 「かぞえる数」と「はかる数」（12ページ、p.14〜25）
%% ---- この章だけで使う絵の部品（頭文字 \bkiib…） ----
\makeatletter
% mm の目もりのあるものさし  \bkiibruler{(x,y)}{1 mm の長さ（図の mm）}{いちばん大きい cm}
%   (x,y) は 0 の目もりの上のはし。目もりは下へのびる。数字は 1 cm ごと
\newcommand\bkiibruler[3]{\begin{scope}[shift={#1}]
  \pgfmathsetmacro\bkiib@w{#3*10*#2}
  \fill[bkYellowBg, draw=bkYellow!80!black, line width=0.8pt, rounded corners=1mm] (-3,0) rectangle (\bkiib@w+3,-11);
  \pgfmathtruncatemacro\bkiib@n{#3*10}
  \foreach \i in {0,...,\bkiib@n} {\pgfmathtruncatemacro\bkiib@c{mod(\i,10)==0}\pgfmathtruncatemacro\bkiib@h{mod(\i,5)==0}
    \ifnum\bkiib@c=1 \draw[bkInk, line width=0.6pt] (\i*#2,0) -- (\i*#2,-4.2);
      \pgfmathtruncatemacro\bkiib@v{\i/10}\node[font=\sffamily\scriptsize, text=bkInk, inner sep=0pt] at (\i*#2,-7.3) {\bkiib@v};
    \else\ifnum\bkiib@h=1 \draw[bkInk, line width=0.4pt] (\i*#2,0) -- (\i*#2,-3);
    \else\draw[bkInk, line width=0.3pt] (\i*#2,0) -- (\i*#2,-1.8);\fi\fi}
\end{scope}}
% 位の表の数字  \bkiibdigit{(x,y)}{数字}
\newcommand\bkiibdigit[2]{\node[font=\sffamily\bfseries, text=bkInk, inner sep=0pt] at #1 {#2};}
% 本だなの正面図（寸法は mm）  \bkiibshelf{実物 100 mm を何 mm でかくか}{はば}{高さ}{たなの段の数}
\newcommand\bkiibshelf[4]{%
  \pgfmathsetmacro\bkiib@sw{#2/100*#1}\pgfmathsetmacro\bkiib@sh{#3/100*#1}\pgfmathsetmacro\bkiib@st{\bkiib@sh/#4}%
  \pgfmathtruncatemacro\bkiib@gap{#3/#4}\pgfmathtruncatemacro\bkiib@n{#4-1}\pgfmathtruncatemacro\bkiib@m{#4-2}%
  \draw[bkInk, line width=0.9pt] (0,0) rectangle (\bkiib@sw,\bkiib@sh);
  \foreach \k in {1,...,\bkiib@n} \draw[bkInk, line width=0.6pt] (0.9,\k*\bkiib@st) -- (\bkiib@sw-0.9,\k*\bkiib@st);
  \draw[bkGray, line width=0.4pt] (0,-1.5) -- (0,-6) (\bkiib@sw,-1.5) -- (\bkiib@sw,-6) (-1.5,0) -- (-6,0) (-1.5,\bkiib@sh) -- (-6,\bkiib@sh);
  \draw[bkBlue, line width=0.5pt, {Stealth[length=3pt]}-{Stealth[length=3pt]}] (0,-4.5) -- (\bkiib@sw,-4.5);
  \node[anchor=north, font=\sffamily\footnotesize, text=bkBlue!80!black, inner sep=0pt] at (\bkiib@sw/2,-6.5) {#2};
  \draw[bkBlue, line width=0.5pt, {Stealth[length=3pt]}-{Stealth[length=3pt]}] (-4.5,0) -- (-4.5,\bkiib@sh);
  \node[anchor=east, font=\sffamily\footnotesize, text=bkBlue!80!black, inner sep=0pt] at (-6.5,\bkiib@sh/2) {#3};
  \draw[bkGray, line width=0.4pt] (\bkiib@sw+1.5,\bkiib@m*\bkiib@st) -- (\bkiib@sw+6,\bkiib@m*\bkiib@st) (\bkiib@sw+1.5,\bkiib@n*\bkiib@st) -- (\bkiib@sw+6,\bkiib@n*\bkiib@st);
  \draw[bkBlue, line width=0.5pt, {Stealth[length=3pt]}-{Stealth[length=3pt]}] (\bkiib@sw+4.5,\bkiib@m*\bkiib@st) -- (\bkiib@sw+4.5,\bkiib@n*\bkiib@st);
  \node[anchor=west, font=\sffamily\footnotesize, text=bkBlue!80!black, inner sep=0pt] at (\bkiib@sw+6.5,\bkiib@m*\bkiib@st+\bkiib@st/2) {\bkiib@gap};
  \node[anchor=north, font=\sffamily\footnotesize, text=bkGray, inner sep=0pt] at (\bkiib@sw/2,-12) {\R{単位}{たん|い}：mm};}
\makeatother

\bkchapter{2}{1}{「かぞえる\R{数}{かず}」と「はかる\R{数}{かず}」}{\bklevel{1}}{「かぞえる\R{数}{かず}」と「はかる\R{数}{かず}」}
\bkmichishirube{1-8,1-11,1-16}{2-3,2-6,2-10}

\begin{goal}
\begin{itemize}[label={\color{bkTheme}●}, leftmargin=1.6\zw, itemsep=0pt, topsep=1pt]
\item \bklevel{1}\enspace「かぞえる\R{数}{かず}」と「はかる\R{数}{かず}」のちがいを\R{説明}{せつ|めい}できる。
\item \bklevel{1}\enspace\R{量}{りょう}を「\R{数値}{すう|ち}×\R{単位}{たん|い}」で\R{表}{あらわ}し、\R{単位}{たん|い}をかえて\R{表}{あらわ}せる。
\item \bklevel{1}\enspace\R{小数}{しょう|すう}の\R{位}{くらい}がわかり、10\R{倍}{ばい}・$\dfrac{1}{10}$ で\R{小数点}{しょう|すう|てん}を\R{動}{うご}かせる。
\end{itemize}
\end{goal}

\begin{hatena}
ソラから、スマホにメッセージがとどきました。「\R{駅}{えき}まで、あと5！」\quad この「5」で、ソラがいつ\R{着}{つ}くか、わかるかな？
\begin{bkfigure}
\begin{tikzpicture}[x=1mm, y=1mm]
  % スマホ
  \fill[bkInk!80, rounded corners=3mm] (0,-1) rectangle (26,30);
  \fill[white, rounded corners=1mm] (2.3,4) rectangle (23.7,26);
  \node[fill=bkGreen!25, rounded corners=2mm, font=\sffamily\small, align=center, inner sep=1.6mm] at (13,15)
    {\R{駅}{えき}まで\\あと5！};
  \fill[bkGray!60] (13,1.5) circle (1.3);
  % 3つの読み方
  \node[font=\sffamily\bfseries\small, text=bkYellow!50!black, inner sep=0pt] at (79,28.5) {どの「5」かな？};
  \foreach \x/\n/\t in {47/a/{5\R{分}{ふん}？\\[2pt]{\footnotesize もうすぐ\R{着}{つ}く}},
                       79/b/{\qty{5}{km}？\\[2pt]{\footnotesize \R{歩}{ある}くと1\R{時間}{じ|かん}\\より\R{長}{なが}い}},
                       111/c/{5\R{駅}{えき}？\\[2pt]{\footnotesize \R{電車}{でん|しゃ}で\\10\R{分}{ぷん}くらい}}} {
    \node[bk node, text width=21mm, minimum height=22mm, inner ysep=1.5mm] (\n) at (\x,12) {\t};}
\end{tikzpicture}
\end{bkfigure}
\end{hatena}

\begin{omoidasou}
\bkrecall{1-8}{\R{長}{なが}さは「\qty{1}{cm} のいくつ\R{分}{ぶん}」で\R{表}{あらわ}す。$\qty{1}{m}=\qty{100}{cm}$、$\qty{1}{cm}=\qty{10}{mm}$}
\bkrecall{1-9}{$\qty{1}{kg}=\qty{1000}{g}$、$\qty{1}{L}=\qty{10}{dL}=\qty{1000}{mL}$}
\bkrecall{1-16}{1を10\R{等分}{とう|ぶん}した1つ\R{分}{ぶん}が0.1。$\qty{1}{mm}=\qty{0.1}{cm}$}
\end{omoidasou}

\clearpage
%% ---------------- p.15 やってみよう・数学の道具 ----------------
\begin{yattemiyou}[\R{筆箱}{ふで|ばこ}の\R{中}{なか}を、かぞえる・はかる]
\begin{enumerate}[label={\bkcircled[bkGreen]{\arabic*}}, leftmargin=2.4\zw, itemsep=1pt, topsep=2pt]
\item \R{筆箱}{ふで|ばこ}の\R{中}{なか}のえんぴつの\R{本数}{ほん|すう}を\R{数}{かぞ}えよう。
\item いちばん\R{長}{なが}いえんぴつの\R{長}{なが}さを、mm の\R{目}{め}もりまで\R{読}{よ}もう。
\item \R{本数}{ほん|すう}と\R{長}{なが}さ。7.5 のような\R{数}{かず}になることがあるのは、どっち？
\end{enumerate}
\begin{bkfigure}
\begin{tikzpicture}[x=1mm, y=1mm]
  \bkiibruler{(0,0)}{0.72}{15}
  \bkpencil{(2.2,3.6)}{100.76}{bkBlue}
  \draw[bkRed, line width=0.8pt] (102.96,6.5) -- (102.96,-1.5);
  \node[anchor=south, font=\sffamily\bfseries\small, text=bkRed!85!black, inner sep=0pt] at (102.96,8) {14\,cm\,3\,mm};
\end{tikzpicture}
\end{bkfigure}
\end{yattemiyou}


%% ---------------- 数学の道具 ----------------
\bksection{doguu}{\R{数学}{すう|がく}の\R{道具}{どう|ぐ}}

\bksubsection{1}{とびとびの\R{数}{かず}と、\R{切}{き}れ\R{目}{め}のない\R{数}{かず}}

\leavevmode\label{bk:2-1:kazoeru}\bkyattemiyouref で、えんぴつは「7\R{本}{ほん}」のように\R{数}{かぞ}えられ、「7.5\R{本}{ほん}」はありません。でも、\R{長}{なが}さは、14\,cm と15\,cm のあいだの14\,cm\,3\,mm でした。

\begin{bkfigure}
\begin{tikzpicture}[x=1mm, y=1mm]
  % 上：本数（とびとび）　下：長さ（切れ目がない）。どちらも 1目もり 11 mm
  \node[anchor=east, font=\sffamily\small, text=bkInk, inner sep=0pt] at (-3,15) {\R{本数}{ほん|すう}};
  \draw[bkLine, line width=0.8pt] (0,15) -- (110,15);
  \foreach \i in {0,...,10} {\fill[bkOrange] (\i*11,15) circle (1.5);
    \node[anchor=north, font=\sffamily\footnotesize, inner sep=0pt] at (\i*11,12.4) {\i};}
  \node[anchor=east, font=\sffamily\small, text=bkInk, inner sep=0pt] at (-3,3) {\R{長}{なが}さ};
  \fill[bkBlue!70] (0,2.2) rectangle (110,3.8);
  \foreach \i in {0,...,10} {\draw[bkInk, line width=0.6pt] (\i*11,1.2) -- (\i*11,4.8);
    \node[anchor=north, font=\sffamily\footnotesize, inner sep=0pt] at (\i*11,-0.2) {\i};}
  % 7 と 8 のあいだを拡大
  \draw[bkGray, line width=0.5pt, densely dashed] (77,-4) -- (40,-9) (88,-4) -- (110,-9);
  \fill[bkBlue!70] (40,-11.3) rectangle (110,-9.7);
  \foreach \k in {0,...,10} {\draw[bkInk, line width=0.5pt] (40+\k*7,-12.3) -- (40+\k*7,-8.7);}
  \foreach \k/\l in {0/7, 3/7.3, 5/7.5, 10/8} {\node[anchor=north, font=\sffamily\footnotesize, inner sep=0pt] at (40+\k*7,-13.7) {\l};}
  \node[anchor=east, font=\sffamily\footnotesize, text=bkBlue!80!black, align=right, inner sep=0pt] at (35,-10.5) {7と8のあいだにも\\\R{数}{かず}がある};
\end{tikzpicture}
\bkfigcaption{\R{本数}{ほん|すう}は\R{点}{てん}で、\R{長}{なが}さは\R{切}{き}れ\R{目}{め}のない\R{線}{せん}で\R{表}{あらわ}せる。}
\end{bkfigure}

\R{人数}{にん|ずう}や\R{本数}{ほん|すう}のように、1つ、2つ、…と\R{数}{かぞ}えて\R{表}{あらわ}す\R{数}{かず}を、この\R{本}{ほん}では\textbf{かぞえる\R{数}{かず}}とよびます。\R{長}{なが}さ・\R{重}{おも}さ・\R{時間}{じ|かん}のように、はかって\R{表}{あらわ}す\R{数}{かず}は\textbf{はかる\R{数}{かず}}です。はかる\R{数}{かず}は\R{切}{き}れ\R{目}{め}がなく、\R{目}{め}もりを\R{細}{こま}かくすれば、\R{小数}{しょう|すう}でいくらでも\R{細}{こま}かく\R{表}{あらわ}せます。\bkidx{かぞえるかず}{かぞえる\R{数}{かず}}\bkidx{はかるかず}{はかる\R{数}{かず}}

\noindent\begin{minipage}[c]{0.60\linewidth}
\R{時間}{じ|かん}でもたしかめてみましょう。50\,m\R{走}{そう}のタイムを、1\R{秒}{びょう}ごとにしか\R{表示}{ひょう|じ}しない\R{時計}{と|けい}ではかると「9\R{秒}{びょう}」。$\dfrac{1}{10}$\R{秒}{びょう}まで\R{表示}{ひょう|じ}するストップウォッチなら「8.7\R{秒}{びょう}」。もっと\R{細}{こま}かいものなら「8.73\R{秒}{びょう}」。はかる\R{道具}{どう|ぐ}の\R{目}{め}もりを\R{細}{こま}かくするほど、\R{小数}{しょう|すう}のけたがふえていきます。
\end{minipage}\hfill
\begin{minipage}[c]{0.37\linewidth}\centering
\begin{tikzpicture}[x=1mm, y=1mm]
  \bkdisplay{(0,26)}{40}{9}
  \bkdisplay{(0,13)}{40}{8.7}
  \bkdisplay{(0,0)}{40}{8.73}
  \node[anchor=west, font=\sffamily\footnotesize, text=bkGray, inner sep=0pt] at (42,29) {\R{秒}{びょう}};
  \node[anchor=west, font=\sffamily\footnotesize, text=bkGray, inner sep=0pt] at (42,16) {\R{秒}{びょう}};
  \node[anchor=west, font=\sffamily\footnotesize, text=bkGray, inner sep=0pt] at (42,3) {\R{秒}{びょう}};
\end{tikzpicture}
\end{minipage}

\par\vspace{1mm}
かぞえる\R{数}{かず}か、はかる\R{数}{かず}かを\R{見分}{み|わ}けるには、「7.5\R{本}{ほん}」「0.5\R{人}{にん}」のような、とちゅうの\R{数}{かず}になることがあるかを\R{考}{かんが}えます。とちゅうの\R{数}{かず}にならないのが、かぞえる\R{数}{かず}です。

\bksubsection{2}{\R{量}{りょう}＝\R{数値}{すう|ち}×\R{単位}{たん|い}}

\leavevmode\label{bk:2-1:tani}\qty{5}{m} のテープを、\qty{1}{m} ずつに\R{切}{き}ってみましょう。

\begin{bkfigure}
\begin{tikzpicture}[x=1mm, y=1mm]
  \foreach \i in {0,...,4} {\bktape{(\i*22,0)}{22}{bkTheme}{\qty{1}{m}}}
  \bkbrace{(0,-1.5)}{(110,-1.5)}{\qty{1}{m} の5つ\R{分}{ぶん}で、\qty{5}{m}}{-1}
\end{tikzpicture}
\end{bkfigure}

\qty{5}{m} は、\qty{1}{m} が5つ\R{分}{ぶん}の\R{長}{なが}さです。もとにする\R{大}{おお}きさ（\qty{1}{m}）を\textbf{\R{単位}{たん|い}}、その\R{何}{なん}こ\R{分}{ぶん}かを\R{表}{あらわ}す\R{数}{かず}（5）を\textbf{\R{数値}{すう|ち}}といいます。\bkidx{たんい}{\R{単位}{たん|い}}\bkidx{すうち}{\R{数値}{すう|ち}}
\[ \qty{5}{m}=5\times\qty{1}{m}\qquad\text{（\R{量}{りょう}＝\R{数値}{すう|ち}×\R{単位}{たん|い}）} \]
なぜ「×」なのでしょう。\qty{1}{m}が5つ\R{分}{ぶん}なので、$\qty{1}{m}+\qty{1}{m}+\qty{1}{m}+\qty{1}{m}+\qty{1}{m}$ です。\R{同}{おな}じ\R{数}{かず}を\R{何回}{なん|かい}もたすことは、かけ\R{算}{ざん}で\R{書}{か}けました（1-12）。だから $\qty{1}{m}\times5$ です。かけ\R{算}{ざん}は\R{入}{い}れかえても\R{答}{こた}えが\R{同}{おな}じなので（1-13）、$5\times\qty{1}{m}$ とも\R{書}{か}けます。\R{物理}{ぶつ|り}では、\R{数値}{すう|ち}を\R{前}{まえ}、\R{単位}{たん|い}を\R{後}{うし}ろに\R{書}{か}くやくそくです。

\noindent\begin{minipage}[c]{0.50\linewidth}\centering
\begin{tikzpicture}[x=1mm, y=1mm]
  \foreach \i in {0,1,2} {\bkmasu{(\i*20,0)}{11}{1}{0}
    \node[anchor=north, font=\sffamily\footnotesize, inner sep=0pt] at (\i*20,-2.5) {\qty{1}{L}};}
  \bkbrace{(-6,-8)}{(46,-8)}{\qty{1}{L}の3つ\R{分}{ぶん}で、\qty{3}{L}}{-1}
\end{tikzpicture}
\end{minipage}\hfill
\begin{minipage}[c]{0.47\linewidth}
かさでも\R{同}{おな}じです。\qty{1}{L}ますで3ばい\R{分}{ぶん}の\R{水}{みず}は、\qty{1}{L}が3つ\R{分}{ぶん}なので、$\qty{3}{L}=3\times\qty{1}{L}$。\R{単位}{たん|い}は\qty{1}{L}、\R{数値}{すう|ち}は3です。
\end{minipage}

\par\vspace{2mm}
\bkhatenaref の「あと5」には、\R{単位}{たん|い}がありませんでした。だから、5\R{分}{ふん}なのか\qty{5}{km}なのか、わからなかったのです。\R{量}{りょう}は、\R{数値}{すう|ち}と\R{単位}{たん|い}がそろって、はじめて\R{伝}{つた}わります。

\bksubsection{3}{\R{単位}{たん|い}をかえると、\R{数値}{すう|ち}がかわる}

\leavevmode\R{同}{おな}じ\R{消}{け}しゴムの\R{長}{なが}さを、\qty{1}{cm} と \qty{1}{mm} で\R{数}{かぞ}えてみましょう。

\begin{bkfigure}
\begin{tikzpicture}[x=1mm, y=1mm]
  % 消しゴム 4 cm（図の上で 110 mm。長さ：高さ＝約 2.3：1）
  \colorlet{bkiibcover}{bkBlue!55}\bkeraser[bkiibcover]{(0,12)}{110}{40}
  \node[anchor=west, font=\sffamily\footnotesize, text=bkGray, align=left, inner sep=0pt] at (112,32) {\R{消}{け}しゴム};
  \draw[bkGray, line width=0.4pt, densely dashed] (0,12) -- (0,-4) (110,12) -- (110,-4);
  % 1 cm ずつ
  \foreach \i in {0,...,3} {\bktape{(\i*27.5,3)}{27.5}{bkOrange}{\qty{1}{cm}}}
  \node[anchor=west, font=\sffamily\small, inner sep=0pt] at (112,6.5) {4こ};
  % 1 mm ずつ
  \foreach \i in {0,...,39} {\filldraw[fill=bkGreen!25, draw=bkGreen!80!black, line width=0.4pt] (\i*2.75,-4) rectangle ++(2.75,5);}
  \node[anchor=west, font=\sffamily\small, inner sep=0pt] at (112,-1.5) {40こ};
  \draw[bkInk!70, line width=0.5pt] (1.4,-5) -- (1.4,-8);
  \node[anchor=west, font=\sffamily\footnotesize, inner sep=0pt] at (3,-8.5) {1つが \qty{1}{mm}};
\end{tikzpicture}
\bkfigcaption{\R{同}{おな}じ\R{長}{なが}さでも、\R{単位}{たん|い}が\R{小}{ちい}さいと\R{数値}{すう|ち}が\R{大}{おお}きくなる。}
\end{bkfigure}

\qty{1}{cm} で\R{数}{かぞ}えると4こ\R{分}{ぶん}、\qty{1}{mm} で\R{数}{かぞ}えると40こ\R{分}{ぶん}。\R{長}{なが}さは\R{同}{おな}じなのに、\R{数値}{すう|ち}は4と40です。\qty{1}{cm} の1こが、\qty{1}{mm} の10こになるからです。\R{式}{しき}で\R{書}{か}くと、
\begin{align*}
\qty{4}{cm}&=4\times\qty{1}{cm} &&\text{（\R{量}{りょう}＝\R{数値}{すう|ち}×\R{単位}{たん|い}）}\\
&=4\times\qty{10}{mm} &&\text{（\qty{1}{cm} を、\R{同}{おな}じ\R{量}{りょう}の \qty{10}{mm} におきかえる）}\\
&=\qty{40}{mm}
\end{align*}
\R{単位}{たん|い}をかえるときは、\R{単位}{たん|い}を\R{同}{おな}じ\R{量}{りょう}におきかえます。

\R{反対}{はん|たい}に、\qty{40}{mm}を cm にするときはどうでしょう。\qty{10}{mm}が\R{集}{あつ}まると\qty{1}{cm}になります。40は10が4つ\R{分}{ぶん}なので、
\begin{align*}
\qty{40}{mm}&=4\times\qty{10}{mm} &&\text{（40は10の4つ\R{分}{ぶん}）}\\
&=4\times\qty{1}{cm} &&\text{（\qty{10}{mm} を、\R{同}{おな}じ\R{量}{りょう}の \qty{1}{cm} におきかえる）}\\
&=\qty{4}{cm}
\end{align*}
\R{同}{おな}じ\R{考}{かんが}え\R{方}{かた}なら、どんな\R{単位}{たん|い}でもかえられます。\R{大}{おお}きい\R{単位}{たん|い}から\R{小}{ちい}さい\R{単位}{たん|い}にすると\R{数値}{すう|ち}は\R{大}{おお}きくなり、\R{小}{ちい}さい\R{単位}{たん|い}から\R{大}{おお}きい\R{単位}{たん|い}にすると\R{数値}{すう|ち}は\R{小}{ちい}さくなります。\R{量}{りょう}（\R{消}{け}しゴムの\R{長}{なが}さ）は、どちらでも\R{同}{おな}じです。

\bksubsection{4}{$\dfrac{1}{10}$の\R{位}{くらい}、$\dfrac{1}{100}$の\R{位}{くらい}}

\leavevmode\label{bk:2-1:kurai}ソラの\R{身長}{しん|ちょう}は \qty{1}{m} と \qty{45}{cm}。m だけで\R{表}{あらわ}すと、\R{何}{なん}m かな？\quad \qty{45}{cm} は \qty{1}{m} より\R{短}{みじか}いので、\qty{1}{m} を\R{細}{こま}かく\R{分}{わ}けて\R{考}{かんが}えます。

\par\vspace{1mm}
\noindent{\sffamily\bfseries\color{bkDeep}①\enspace\qty{1}{m} を10\R{等分}{とう|ぶん}する}\par
\qty{1}{m}（\qty{100}{cm}）を10\R{等分}{とう|ぶん}すると、1つ\R{分}{ぶん}は $100\div10=10$ で \qty{10}{cm}。1を10\R{等分}{とう|ぶん}した1つ\R{分}{ぶん}は0.1でした（1-16）。だから、$\qty{10}{cm}=\qty{0.1}{m}$ です。

\par\vspace{1mm}
\noindent{\sffamily\bfseries\color{bkDeep}②\enspace\qty{0.1}{m} を、さらに10\R{等分}{とう|ぶん}する}\par
\qty{10}{cm} をさらに10\R{等分}{とう|ぶん}すると、1つ\R{分}{ぶん}は $10\div10=1$ で \qty{1}{cm}。0.1を10\R{等分}{とう|ぶん}した1つ\R{分}{ぶん}の\R{数}{かず}を \textbf{0.01} と\R{書}{か}き、「れい\R{点}{てん}れい\R{一}{いち}」と\R{読}{よ}みます。だから、$\qty{1}{cm}=\qty{0.01}{m}$ です。

\noindent\begin{minipage}[c]{0.50\linewidth}\centering
\begin{tikzpicture}[x=1mm, y=1mm]
  % 大きな正方形を 1 とする（1ます 4 mm）
  \fill[bkOrange!30] (0,0) rectangle (4,40);
  \fill[bkBlue!55] (36,0) rectangle (40,4);
  \draw[bkGray!70, line width=0.3pt] (0,0) grid[step=4] (40,40);
  \draw[bkInk, line width=0.9pt] (0,0) rectangle (40,40);
  \draw[bkInk!70, line width=0.5pt] (40.8,38) -- (44.5,38);
  \node[anchor=west, font=\sffamily\footnotesize, inner sep=0pt] at (45.5,38) {\R{全体}{ぜん|たい}：1};
  \draw[bkInk!70, line width=0.5pt] (2,40.8) -- (2,44);
  \node[anchor=south west, font=\sffamily\footnotesize, text=bkOrange!80!black, inner sep=0pt] at (0,45.3) {たて1\R{列}{れつ}：0.1};
  \draw[bkInk!70, line width=0.5pt] (40.8,2) -- (44.5,2);
  \node[anchor=west, font=\sffamily\footnotesize, text=bkBlue!80!black, inner sep=0pt] at (45.5,2) {1ます：0.01};
\end{tikzpicture}
\end{minipage}\hfill
\begin{minipage}[c]{0.47\linewidth}
\R{大}{おお}きな\R{正方形}{せい|ほう|けい}を1とします。たてに10\R{列}{れつ}に\R{分}{わ}けた1\R{列}{れつ}が0.1、それをさらに10に\R{分}{わ}けた1ますが0.01です。\\[2pt]
0.01が10こで0.1、0.1が10こで1。だから、0.01が100こで1です。0.01は、1を100\R{等分}{とう|ぶん}した1つ\R{分}{ぶん}でもあります（\qty{1}{m} の\R{中}{なか}に \qty{1}{cm} は100こ）。
\end{minipage}

\par\vspace{2mm}
\noindent{\sffamily\bfseries\color{bkDeep}③\enspace\qty{45}{cm} を m で\R{表}{あらわ}す}\par
$\qty{45}{cm}=\qty{40}{cm}+\qty{5}{cm}$ と\R{分}{わ}けます。\qty{40}{cm} は \qty{10}{cm}（\qty{0.1}{m}）が4こで \qty{0.4}{m}。\qty{5}{cm} は \qty{1}{cm}（\qty{0.01}{m}）が5こで \qty{0.05}{m}。\R{合}{あ}わせて \qty{0.45}{m} です。ソラの\R{身長}{しん|ちょう}は、$\qty{1}{m}+\qty{0.45}{m}$ で \textbf{\qty{1.45}{m}}（いってんよんご メートル）です。

\begin{bkfigure}
\begin{tikzpicture}[x=1mm, y=1mm]
  % 1.45 を正方形で（1ます 3.6 mm）
  \fill[bkTheme!35] (0,0) rectangle (36,36);
  \draw[bkGray!70, line width=0.3pt] (0,0) grid[step=3.6] (36,36);
  \draw[bkInk, line width=0.9pt] (0,0) rectangle (36,36);
  \begin{scope}[shift={(48,0)}]
    \fill[bkOrange!30] (0,0) rectangle (14.4,36);
    \fill[bkBlue!55] (14.4,18) rectangle (18,36);
    \draw[bkGray!70, line width=0.3pt] (0,0) grid[step=3.6] (36,36);
    \draw[bkInk, line width=0.9pt] (0,0) rectangle (36,36);
  \end{scope}
  \node[anchor=north, font=\sffamily\footnotesize, inner sep=0pt] at (18,-2.5) {1が1こ：1};
  \node[anchor=north, font=\sffamily\footnotesize, text=bkOrange!80!black, inner sep=0pt] at (66,-2.5) {0.1が4こ：0.4};
  \node[anchor=north, font=\sffamily\footnotesize, text=bkBlue!80!black, inner sep=0pt] at (66,-7.5) {0.01が5こ：0.05};
  \node[anchor=west, font=\sffamily\bfseries\small, inner sep=0pt] at (89,18) {$1+0.4+0.05=1.45$};
\end{tikzpicture}
\bkfigcaption{1.45 は、1が1こ、0.1が4こ、0.01が5こ。}
\end{bkfigure}

\R{数直線}{すう|ちょく|せん}でも\R{見}{み}てみましょう。1と2のあいだを10\R{等分}{とう|ぶん}した1\R{目}{め}もりが0.1、1.4と1.5のあいだをさらに10\R{等分}{とう|ぶん}した1\R{目}{め}もりが0.01です。

\begin{bkfigure}
\begin{tikzpicture}[x=1mm, y=1mm]
  % 1 m〜2 m（0.1 m ごと）
  \draw[bkInk, line width=0.8pt] (0,11) -- (120,11);
  \foreach \k in {0,...,10} {\draw[bkInk, line width=0.6pt] (\k*12,9.8) -- (\k*12,12.2);}
  \foreach \k/\l in {0/1, 4/1.4, 5/1.5, 10/2} {\node[anchor=south, font=\sffamily\footnotesize, inner sep=0pt] at (\k*12,13.6) {\l};}
  \node[anchor=west, font=\sffamily\footnotesize, text=bkGray, inner sep=0pt] at (122,11) {m};
  \fill[bkBlue!25] (48,9.8) rectangle (60,12.2);
  % 1.4 m〜1.5 m（0.01 m ごと）
  \draw[bkGray, line width=0.5pt, densely dashed] (48,9) -- (0,1) (60,9) -- (120,1);
  \draw[bkInk, line width=0.8pt] (0,-0.5) -- (120,-0.5);
  \foreach \k in {0,...,10} {\draw[bkInk, line width=0.5pt] (\k*12,-1.7) -- (\k*12,0.7);}
  \foreach \k/\l in {0/1.4, 10/1.5} {\node[anchor=north, font=\sffamily\footnotesize, inner sep=0pt] at (\k*12,-3.1) {\l};}
  \draw[bk arrow, bkOrange] (60,-8.5) -- (60,-2.4);
  \node[anchor=west, font=\sffamily\bfseries\small, text=bkOrange!80!black, inner sep=0pt] at (62.5,-7) {1.45（ソラの\R{身長}{しん|ちょう}）};
  \node[anchor=west, font=\sffamily\footnotesize, text=bkGray, inner sep=0pt] at (122,-0.5) {m};
\end{tikzpicture}
\bkfigcaption{\R{上}{うえ}は1\R{目}{め}もり \qty{0.1}{m}、\R{下}{した}は1\R{目}{め}もり \qty{0.01}{m}}
\end{bkfigure}

\noindent\begin{minipage}[c]{0.50\linewidth}
\begin{bktable}{colspec={*{3}{Q[c,m,wd=14mm]}}, rowsep=1pt, colsep=3pt, cells={font=\footnotesize}, row{1}={font=\sffamily\bfseries\footnotesize}, row{2}={font=\large}}
\R{一}{いち}の\R{位}{くらい} & $\frac{1}{10}$の\R{位}{くらい} & $\frac{1}{100}$の\R{位}{くらい} \\
1 & 4 & 5 \\
{\qty{1}{m}\\が1こ} & {\qty{0.1}{m}\\が4こ} & {\qty{0.01}{m}\\が5こ} \\
\end{bktable}
\end{minipage}\hfill
\begin{minipage}[c]{0.48\linewidth}
\qty{1}{m} が1こ、\qty{0.1}{m} が4こ、\qty{0.01}{m} が5こで、\textbf{\qty{1.45}{m}}。\R{数字}{すう|じ}を\R{書}{か}く\R{場所}{ば|しょ}（\R{位}{くらい}）で、0.1が\R{何}{なん}こか、0.01が\R{何}{なん}こかを\R{表}{あらわ}しています。
\end{minipage}

\par\vspace{2mm}
\R{小数点}{しょう|すう|てん}のすぐ\R{右}{みぎ}を$\dfrac{1}{10}$の\R{位}{くらい}（\R{小数第}{しょう|すう|だい}1\R{位}{い}）、その\R{右}{みぎ}を$\dfrac{1}{100}$の\R{位}{くらい}（\R{小数第}{しょう|すう|だい}2\R{位}{い}）といいます。\bkidx{しょうすうのくらい}{\R{小数}{しょう|すう}の\R{位}{くらい}}$\dfrac{1}{10}$は「10\R{分}{ぶん}の1」と\R{読}{よ}み、1を10\R{等分}{とう|ぶん}した1つ\R{分}{ぶん}（0.1）のことです。$\dfrac{1}{100}$は「100\R{分}{ぶん}の1」で、0.01のことです。

\par\vspace{1mm}
\noindent{\sffamily\bfseries\color{bkDeep}1.45は、0.01を\R{何}{なん}こ\R{集}{あつ}めた\R{数}{かず}？}\par
1は0.01が100こ。0.4は0.1が4こで、0.1は0.01が10こなので、0.01が $4\times10=40$ こ。0.05は0.01が5こ。\R{合}{あ}わせて $100+40+5=145$ で、145こです。\R{単位}{たん|い}で\R{考}{かんが}えても\R{同}{おな}じです。$\qty{1.45}{m}=\qty{145}{cm}$ で、$\qty{1}{cm}=\qty{0.01}{m}$ なので、\qty{0.01}{m} が145こです。

\bksubsection{5}{10\R{倍}{ばい}すると、\R{小数点}{しょう|すう|てん}が\R{右}{みぎ}へ1つ}

\leavevmode\label{bk:2-1:ten}\qty{1.45}{m} を cm で\R{表}{あらわ}しましょう。\qty{1}{m} は \qty{100}{cm} なので、\qty{1.45}{m} は \qty{100}{cm} の1.45こ\R{分}{ぶん}です。$1.45\times100$ を\R{求}{もと}めればよいのです。まず、10\R{倍}{ばい}するとどうなるかを\R{考}{かんが}えます。

\par\vspace{1mm}
\noindent{\sffamily\bfseries\color{bkDeep}①\enspace 10\R{倍}{ばい}すると、なぜ\R{位}{くらい}が1つ\R{上}{あ}がるの？}\par
\qty{1.45}{m} のテープを10\R{本}{ぽん}つなぐと\R{考}{かんが}えます。\R{位}{くらい}ごとに\R{分}{わ}けて、それぞれを10\R{倍}{ばい}します。

\begin{bkcenter}
\begin{bktable}{colspec={Q[l,m,wd=26mm] Q[c,m,wd=14mm] Q[l,m,wd=70mm]}, rowsep=1pt, cells={font=\small}, row{1}={font=\sffamily\bfseries\small}}
1\R{本}{ぽん}の\R{中}{なか}み & 10\R{本}{ぽん}\R{分}{ぶん} & わけ \\
\qty{1}{m} & \qty{10}{m} & 1が10こで10。\R{一}{いち}の\R{位}{くらい}の1が、\R{十}{じゅう}の\R{位}{くらい}へ。 \\
\qty{0.4}{m}（0.1が4こ） & \qty{4}{m} & 0.1が40こ。0.1が10こで1なので、4。 \\
\qty{0.05}{m}（0.01が5こ） & \qty{0.5}{m} & 0.01が50こ。0.01が10こで0.1なので、0.5。 \\
\end{bktable}
\bktabcaption{10こ\R{集}{あつ}まると、1つ\R{上}{うえ}の\R{位}{くらい}の1こになる。}
\end{bkcenter}

\R{合}{あ}わせて $10+4+0.5=14.5$ で、\qty{14.5}{m}。どの\R{位}{くらい}の\R{数}{かず}も、10こ\R{集}{あつ}まると1つ\R{上}{うえ}の\R{位}{くらい}の1こになります（\pageref{bk:2-1:kurai}ページ）。だから、10\R{倍}{ばい}すると、どの\R{数字}{すう|じ}も\R{位}{くらい}が1つ\R{上}{あ}がるのです。

\noindent\begin{minipage}[c]{0.52\linewidth}
\begin{tikzpicture}[x=1mm, y=1mm]
  \foreach \c/\t in {0/\R{百}{ひゃく}, 1/\R{十}{じゅう}, 2/\R{一}{いち}, 3/$\frac{1}{10}$, 4/$\frac{1}{100}$} {
    \node[font=\sffamily\footnotesize, text=bkGray, inner sep=0pt] at (\c*11,7.6) {\t};}
  \foreach \r in {0,1,2} {\foreach \c in {0,...,4} {\draw[bkLine, line width=0.6pt] (\c*11-5,-\r*10-4.2) rectangle ++(10,8.4);}
    \fill[bkRed] (27.5,-\r*10-3) circle (0.9);}
  \bkiibdigit{(22,0)}{1}\bkiibdigit{(33,0)}{4}\bkiibdigit{(44,0)}{5}
  \bkiibdigit{(11,-10)}{1}\bkiibdigit{(22,-10)}{4}\bkiibdigit{(33,-10)}{5}
  \bkiibdigit{(0,-20)}{1}\bkiibdigit{(11,-20)}{4}\bkiibdigit{(22,-20)}{5}
  \draw[bk arrow, bkOrange, line width=1pt] (51,0) to[bend left=45] node[right, font=\sffamily\small, text=bkOrange!80!black] {×10} (51,-10);
  \draw[bk arrow, bkOrange, line width=1pt] (51,-10) to[bend left=45] node[right, font=\sffamily\small, text=bkOrange!80!black] {×10} (51,-20);
\end{tikzpicture}
\end{minipage}\hfill
\begin{minipage}[c]{0.46\linewidth}
10\R{倍}{ばい}すると、どの\R{数字}{すう|じ}も\R{位}{くらい}が1つ\R{上}{あ}がります。\R{数字}{すう|じ}のならび（1，4，5）はそのままで、\R{小数点}{しょう|すう|てん}（\R{赤}{あか}い\R{点}{てん}）が\R{右}{みぎ}へ1つ\R{動}{うご}いたように\R{見}{み}えます。
\end{minipage}

\par\vspace{2mm}
\noindent{\sffamily\bfseries\color{bkDeep}②\enspace 100\R{倍}{ばい}、1000\R{倍}{ばい}}\par
$100=10\times10$ なので、100\R{倍}{ばい}は10\R{倍}{ばい}を2\R{回}{かい}することです。\allowbreak\mbox{1.45 → 14.5}\allowbreak\mbox{ → 145} なので、$\qty{1.45}{m}=1.45\times\qty{100}{cm}=\qty{145}{cm}$。$1000=10\times10\times10$ なので、1000\R{倍}{ばい}は10\R{倍}{ばい}を3\R{回}{かい}。\R{小数点}{しょう|すう|てん}は\R{右}{みぎ}へ3つ\R{動}{うご}きます。

\noindent\begin{minipage}[c]{0.62\linewidth}
\R{電卓}{でん|たく}で\mbox{「1.45 × 10 ＝」}とおすと 14.5、つづけて\mbox{「× 10 ＝」}とおすと 145。\R{数字}{すう|じ}のならびはそのままで、\R{小数点}{しょう|すう|てん}の\R{場所}{ば|しょ}だけがかわります。
\end{minipage}\hfill
\begin{minipage}[c]{0.35\linewidth}\centering
\begin{tikzpicture}[x=1mm, y=1mm]
  \bkdisplay{(0,14)}{38}{14.5}
  \bkdisplay{(0,0)}{38}{145}
  \node[anchor=east, font=\sffamily\footnotesize, text=bkOrange!80!black, inner sep=0pt] at (-1.5,19.5) {×10};
  \node[anchor=east, font=\sffamily\footnotesize, text=bkOrange!80!black, inner sep=0pt] at (-1.5,5.5) {×10};
\end{tikzpicture}
\end{minipage}

\par\vspace{2mm}
\noindent{\sffamily\bfseries\color{bkDeep}③\enspace $\dfrac{1}{10}$にする（10でわる）と、\R{位}{くらい}が1つ\R{下}{さ}がる}\par
\R{反対}{はん|たい}に、$\dfrac{1}{10}$にする（10でわる）と、どの\R{数字}{すう|じ}も\R{位}{くらい}が1つ\R{下}{さ}がり、\R{小数点}{しょう|すう|てん}が\R{左}{ひだり}へ1つ\R{動}{うご}きます。1を10\R{等分}{とう|ぶん}すると0.1、0.1を10\R{等分}{とう|ぶん}すると0.01だったからです。\qty{145}{cm} を m にするときは、\qty{100}{cm} が \qty{1}{m} なので、100でわります。\mbox{145 → 14.5 → 1.45} で、\qty{1.45}{m}。もとにもどりました。

\par\vspace{1mm}
\noindent{\sffamily\bfseries\color{bkDeep}④\enspace\R{数字}{すう|じ}のない\R{位}{くらい}には0を\R{書}{か}く}\par
0.4を100\R{倍}{ばい}すると、\mbox{0.4 → 4 → 40}。\R{一}{いち}の\R{位}{くらい}に\R{数字}{すう|じ}がなくなるので、0を\R{書}{か}きます。35を100でわると、\mbox{35 → 3.5 → 0.35}。\R{一}{いち}の\R{位}{くらい}があくので、0を\R{書}{か}きます。

\begin{bkfigure}
\begin{tikzpicture}[x=0.8mm, y=0.8mm]
  \useasboundingbox (-5,-24) rectangle (140,10);
  \foreach \ox/\op in {0/×10, 74/÷10} {
    \foreach \c/\t in {0/\R{百}{ひゃく}, 1/\R{十}{じゅう}, 2/\R{一}{いち}, 3/$\frac{1}{10}$, 4/$\frac{1}{100}$} {
      \node[font=\sffamily\footnotesize, text=bkGray, inner sep=0pt] at (\ox+\c*11,7.6) {\t};}
    \foreach \r in {0,1,2} {\foreach \c in {0,...,4} {\draw[bkLine, line width=0.6pt] (\ox+\c*11-5,-\r*10-4.2) rectangle ++(10,8.4);}
      \fill[bkRed] (\ox+27.5,-\r*10-3) circle (0.9);}
    \draw[bk arrow, bkOrange, line width=1pt] (\ox+51,0) to[bend left=45] node[right=1.2mm, font=\sffamily\small, text=bkOrange!80!black] {\op} (\ox+51,-10);
    \draw[bk arrow, bkOrange, line width=1pt] (\ox+51,-10) to[bend left=45] node[right=1.2mm, font=\sffamily\small, text=bkOrange!80!black] {\op} (\ox+51,-20);}
  % 0.4 → 4 → 40
  \bkiibdigit{(22,0)}{0}\bkiibdigit{(33,0)}{4}
  \bkiibdigit{(22,-10)}{4}
  \bkiibdigit{(11,-20)}{4}\node[font=\sffamily\bfseries, text=bkRed, inner sep=0pt] at (22,-20) {0};
  % 35 → 3.5 → 0.35
  \bkiibdigit{(85,0)}{3}\bkiibdigit{(96,0)}{5}
  \bkiibdigit{(96,-10)}{3}\bkiibdigit{(107,-10)}{5}
  \node[font=\sffamily\bfseries, text=bkRed, inner sep=0pt] at (96,-20) {0};\bkiibdigit{(107,-20)}{3}\bkiibdigit{(118,-20)}{5}
\end{tikzpicture}
\bkfigcaption{\R{赤}{あか}い0は、あいた\R{位}{くらい}に\R{書}{か}きたした0。}
\end{bkfigure}

\begin{motto}[1]
\textbf{かぞえる\R{数}{かず}なのに、\R{小数}{しょう|すう}？}\par
\R{日本}{に|ほん}の1つの\R{家}{いえ}に\R{住}{す}む\R{人}{ひと}の\R{数}{かず}は、\R{平均}{へい|きん}すると\R{約}{やく}2.2\R{人}{にん}です（2020\R{年}{ねん}）。たくさんの\R{家}{いえ}の\R{人数}{にん|ずう}をならした\R{数}{かず}（\R{平均}{へい|きん}、2-6）なので、\R{小数}{しょう|すう}になります。2.2\R{人}{にん}の\R{家}{いえ}があるわけではありません。

\par\vspace{1.5mm}
\textbf{「\R{重}{おも}さ」と「\R{質量}{しつ|りょう}」}\par
この\R{章}{しょう}では、kg や g ではかる\R{量}{りょう}を、ふだんのことばで「\R{重}{おも}さ」とよびました。\R{理科}{り|か}や\R{物理}{ぶつ|り}では、この\R{量}{りょう}を「\R{質量}{しつ|りょう}」といいます（2-10）。\R{質量}{しつ|りょう}は、\R{月}{つき}へ\R{持}{も}っていっても\R{変}{か}わりません。
\end{motto}

\begin{taisetsu}
\R{量}{りょう}＝\R{数値}{すう|ち}×\R{単位}{たん|い}。\R{量}{りょう}は、\R{数値}{すう|ち}と\R{単位}{たん|い}をそろえて\R{書}{か}く。\\
\R{単位}{たん|い}をかえるときは、\R{単位}{たん|い}を\R{同}{おな}じ\R{量}{りょう}におきかえる。\\
10\R{倍}{ばい}すると\R{小数点}{しょう|すう|てん}は\R{右}{みぎ}へ1つ、$\dfrac{1}{10}$にすると\R{左}{ひだり}へ1つ\R{動}{うご}く。
\end{taisetsu}

\begin{kotoba}
\small
\bkword{かぞえる\R{数}{かず}}{1つ、2つ、…と\R{数}{かぞ}えて\R{表}{あらわ}す\R{数}{かず}。\R{人数}{にん|ずう}・\R{本数}{ほん|すう}など。とびとびの\R{数}{かず}。}
\bkword{はかる\R{数}{かず}}{\R{長}{なが}さ・\R{重}{おも}さ・\R{時間}{じ|かん}など、はかって\R{表}{あらわ}す\R{数}{かず}。\R{切}{き}れ\R{目}{め}がない。}
\bkword{\R{単位}{たん|い}}{\R{量}{りょう}をはかるときの、もとにする\R{大}{おお}きさ。m、kg、L、\R{秒}{びょう}など。}
\bkword{\R{数値}{すう|ち}}{\R{量}{りょう}が\R{単位}{たん|い}の\R{何}{なん}こ\R{分}{ぶん}かを\R{表}{あらわ}す\R{数}{かず}。}
\end{kotoba}

%% ---------------- 例題 ----------------
\bksection{reidai}{\R{例題}{れい|だい}でたしかめよう}

\begin{reidai}[かぞえる\R{数}{かず}か、はかる\R{数}{かず}か]
\R{次}{つぎ}の\R{量}{りょう}は、かぞえる\R{数}{かず}・はかる\R{数}{かず}のどちらで\R{表}{あらわ}すでしょう。\\
\mbox{（1）ノートの\R{冊数}{さっ|すう}}\quad \mbox{（2）お\R{茶}{ちゃ}の\R{量}{りょう}}\quad \mbox{（3）50\,m\R{走}{そう}のタイム}\quad \mbox{（4）クラスの\R{人数}{にん|ずう}}\quad \mbox{（5）\R{体温}{たい|おん}}\par
\kangaekata 0.5 のような、とちゅうの\R{数}{かず}になることがあるかを\R{考}{かんが}えます。\par
\shiki （1）「2.5\R{冊}{さつ}」のノートはない。\quad（2）「\qty{0.5}{L}」のお\R{茶}{ちゃ}はある。\par
\hspace*{2.6\zw}（3）「8.7\R{秒}{びょう}」のタイムはある。\quad（4）「30.5\R{人}{にん}」のクラスはない。\par
\hspace*{2.6\zw}（5）「\qty{36.5}{\degreeCelsius}」の\R{体温}{たい|おん}はある。\par
\kotae （1）かぞえる\R{数}{かず}\quad（2）はかる\R{数}{かず}\quad（3）はかる\R{数}{かず}\quad（4）かぞえる\R{数}{かず}\quad（5）はかる\R{数}{かず}
\end{reidai}

\begin{jibun}
\R{次}{つぎ}の\R{量}{りょう}は、かぞえる\R{数}{かず}・はかる\R{数}{かず}のどちらで\R{表}{あらわ}すでしょう。\\
\mbox{（1）\R{教室}{きょう|しつ}のいすの\R{数}{かず}}\quad \mbox{（2）ランドセルの\R{重}{おも}さ}\quad \mbox{（3）\R{金魚}{きん|ぎょ}の\R{数}{かず}}\quad \mbox{（4）プールの\R{深}{ふか}さ}\quad \mbox{（5）\R{家族}{か|ぞく}の\R{人数}{にん|ずう}}
\bkjans{（1）\nobreak かぞえる\R{数}{かず}\quad（2）\nobreak はかる\R{数}{かず}\quad（3）\nobreak かぞえる\R{数}{かず}\quad（4）\nobreak はかる\R{数}{かず}\quad（5）\nobreak かぞえる\R{数}{かず}}
\end{jibun}

\begin{reidai}[\R{単位}{たん|い}をかえる]
（1）\qty{3}{m} は\R{何}{なん}cm？\quad（2）\qty{4000}{g} は\R{何}{なん}kg？\\
（3）\qty{3}{m} のテープを \qty{1}{cm} ずつに\R{切}{き}ると\R{何本}{なん|ぼん}？\quad \qty{1}{mm} ずつなら\R{何本}{なん|ぼん}？\par
\kangaekata （1）は \qty{1}{m} を、\R{同}{おな}じ\R{量}{りょう}の \qty{100}{cm} におきかえます。（2）は、4000が1000の4つ\R{分}{ぶん}なので、まず $4\times\qty{1000}{g}$ と\R{書}{か}き、\qty{1000}{g} を \qty{1}{kg} におきかえます。（3）の\R{本数}{ほん|すう}は「\qty{3}{m} が \qty{1}{cm} の\R{何}{なん}こ\R{分}{ぶん}か」なので、cm で\R{表}{あらわ}したときの\R{数値}{すう|ち}です。\par
\shiki （1）$\qty{3}{m}=3\times\qty{1}{m}=3\times\qty{100}{cm}=\qty{300}{cm}$\par
\hspace*{2.6\zw}（2）$\qty{4000}{g}=4\times\qty{1000}{g}=4\times\qty{1}{kg}=\qty{4}{kg}$\par
\hspace*{2.6\zw}（3）（1）から \qty{3}{m} は \qty{300}{cm}。\qty{1}{cm} が300こ\R{分}{ぶん}。\par
\hspace*{4.6\zw}$\qty{300}{cm}=300\times\qty{10}{mm}=\qty{3000}{mm}$。\qty{1}{mm} が3000こ\R{分}{ぶん}。\par
\kotae （1）\qty{300}{cm}\quad（2）\qty{4}{kg}\quad（3）300\R{本}{ぽん}、3000\R{本}{ぼん}
\end{reidai}

\begin{jibun}
（1）\qty{6}{L} は\R{何}{なん}mL？\quad（2）\qty{500}{cm} は\R{何}{なん}m？\\
（3）\qty{4}{m} のひもを \qty{1}{cm} ずつに\R{切}{き}ると\R{何本}{なん|ぼん}？\quad \qty{1}{mm} ずつなら\R{何本}{なん|ぼん}？
\bkjans{（1）\nobreak \qty{6000}{mL}\quad（2）\nobreak \qty{5}{m}\quad（3）\nobreak 400\R{本}{ぽん}、4000\R{本}{ぼん}}
\end{jibun}

\begin{reidai}[cm と mm がまざった\R{長}{なが}さ]
\leavevmode\label{bk:2-1:enpitsu}\bkyattemiyouref のえんぴつの\R{長}{なが}さ 14\,cm\,3\,mm を、（1）cm だけ、（2）mm だけで\R{表}{あらわ}しましょう。\par
\kangaekata 14\,cm\,3\,mm は、\qty{14}{cm} と \qty{3}{mm} をつないだ\R{長}{なが}さです。どちらかの\R{単位}{たん|い}にそろえてから、たします。\qty{1}{mm} は \qty{0.1}{cm}（1-16）、\qty{1}{cm} は \qty{10}{mm} です。\par
\begin{bkcenter}
\begin{tikzpicture}[x=1mm, y=1mm]
  \bktape{(0,0)}{98}{bkBlue}{\qty{14}{cm}}
  \bktape{(98,0)}{2.1}{bkOrange}{}
  \node[anchor=south, font=\sffamily\footnotesize, text=bkOrange!80!black, inner sep=0pt] at (99,8.5) {\qty{3}{mm}};
\end{tikzpicture}
\end{bkcenter}
\shiki （1）$\qty{3}{mm}=3\times\qty{0.1}{cm}=\qty{0.3}{cm}$ なので、\par
\hspace*{4.6\zw}$\qty{14}{cm}+\qty{0.3}{cm}=\qty{14.3}{cm}$\par
\hspace*{2.6\zw}（2）$\qty{14}{cm}=14\times\qty{10}{mm}=\qty{140}{mm}$ なので、\par
\hspace*{4.6\zw}$\qty{140}{mm}+\qty{3}{mm}=\qty{143}{mm}$\par
\kotae （1）\qty{14.3}{cm}\quad（2）\qty{143}{mm}
\end{reidai}

\begin{jibun}
（1）9\,cm\,4\,mm を、mm だけと cm だけで\R{表}{あらわ}しましょう。\quad（2）2\,m\,35\,cm を、cm だけと m だけで\R{表}{あらわ}しましょう。
\bkjans{（1）\nobreak \qty{94}{mm}、\qty{9.4}{cm}\quad（2）\nobreak \qty{235}{cm}、\qty{2.35}{m}}
\end{jibun}

\begin{reidai}[ものさしの\R{目}{め}もりを\R{読}{よ}む]
\R{図}{ず}のクレヨンの\R{長}{なが}さを、（1）○cm○mm、（2）cm だけ、（3）mm だけ で\R{表}{あらわ}しましょう。\par
\begin{bkcenter}
\begin{tikzpicture}[x=1mm, y=1mm]
  \bkiibruler{(0,0)}{0.9}{10}
  \bkpencil{(2.2,3.6)}{68}{bkRed}
  \draw[bkRed, line width=0.8pt] (70.2,6.5) -- (70.2,-1.5);
\end{tikzpicture}
\end{bkcenter}
\kangaekata 0の\R{目}{め}もりから、はしまでを\R{読}{よ}みます。\R{数字}{すう|じ}のついた\R{長}{なが}い\R{目}{め}もりが1\,cm ごと、\R{短}{みじか}い\R{目}{め}もり1つが1\,mm です。\R{中}{なか}くらいの\R{目}{め}もりは5\,mm です。\par
\shiki （1）7\,cm の\R{目}{め}もりをこえて、\R{小}{ちい}さい\R{目}{め}もりが8つ。7\,cm\,8\,mm\par
\hspace*{2.6\zw}（2）$\qty{8}{mm}=\qty{0.8}{cm}$ なので、$\qty{7}{cm}+\qty{0.8}{cm}=\qty{7.8}{cm}$\par
\hspace*{2.6\zw}（3）$\qty{7}{cm}=\qty{70}{mm}$ なので、$\qty{70}{mm}+\qty{8}{mm}=\qty{78}{mm}$\par
\kotae （1）7\,cm\,8\,mm\quad（2）\qty{7.8}{cm}\quad（3）\qty{78}{mm}
\end{reidai}

\begin{jibun}
\R{図}{ず}の\R{消}{け}しゴムの\R{長}{なが}さを、（1）○cm○mm、（2）cm だけ、（3）mm だけ で\R{表}{あらわ}しましょう。\par
\begin{bkcenter}
\begin{tikzpicture}[x=1mm, y=1mm]
  \bkiibruler{(0,0)}{0.9}{8}
  \colorlet{bkiibcover2}{bkGreen!55}\bkeraser[bkiibcover2]{(0,3.6)}{38.7}{16.8}
  \draw[bkRed, line width=0.8pt] (38.7,21.5) -- (38.7,-1.5);
\end{tikzpicture}
\end{bkcenter}
\bkjans{（1）\nobreak 4\,cm\,3\,mm\quad（2）\nobreak \qty{4.3}{cm}\quad（3）\nobreak \qty{43}{mm}}
\end{jibun}

\begin{reidai}[\R{小数}{しょう|すう}の\R{位}{くらい}と10\R{倍}{ばい}]
\qty{2.53}{m} について\R{答}{こた}えましょう。（1）5 は\R{何}{なん}の\R{位}{くらい}の\R{数字}{すう|じ}で、\R{何}{なに}が\R{何}{なん}こあることを\R{表}{あらわ}していますか。（2）2.53 は 0.01 を\R{何}{なん}こ\R{集}{あつ}めた\R{数}{かず}ですか。（3）2.53 を10\R{倍}{ばい}、100\R{倍}{ばい}した\R{数}{かず}と、$\dfrac{1}{10}$にした\R{数}{かず}を\R{求}{もと}めましょう。\par
\kangaekata \R{位}{くらい}ごとに\R{分}{わ}けると、2 は \qty{1}{m} が2こ、5 は \qty{0.1}{m} が5こ、3 は \qty{0.01}{m} が3こ。10\R{倍}{ばい}で\R{小数点}{しょう|すう|てん}は\R{右}{みぎ}へ1つ、$\dfrac{1}{10}$で\R{左}{ひだり}へ1つ\R{動}{うご}きます。\par
\shiki （1）\R{小数点}{しょう|すう|てん}のすぐ\R{右}{みぎ}の\R{数字}{すう|じ}なので、$\dfrac{1}{10}$の\R{位}{くらい}。\par
\hspace*{2.6\zw}（2）1は0.01の100こなので、2は200こ。0.1は0.01の10こなので、\par
\hspace*{4.6\zw}0.5は50こ。0.03は3こ。$200+50+3=253$\par
\hspace*{2.6\zw}（3）$2.53\times10=25.3$\quad（\R{右}{みぎ}へ1つ）\par
\hspace*{4.6\zw}$2.53\times100=253$\quad（\R{右}{みぎ}へ2つ）\par
\hspace*{4.6\zw}$2.53\div10=0.253$\quad（\R{左}{ひだり}へ1つ。\R{一}{いち}の\R{位}{くらい}に0を\R{書}{か}く）\par
\kotae （1）$\dfrac{1}{10}$の\R{位}{くらい}。\qty{0.1}{m} が5こ\quad（2）253こ\quad（3）25.3、253、0.253
\end{reidai}

\begin{jibun}
\qty{4.07}{kg} について、\bkreidairef{5}の（1）〜（3）と\R{同}{おな}じことに\R{答}{こた}えましょう。（1）は 7 について\R{答}{こた}えます。（$\dfrac{1}{10}$の\R{位}{くらい}の0は、\qty{0.1}{kg} が0こ、ということです。）
\bkjans{（1）\nobreak $\dfrac{1}{100}$の\R{位}{くらい}。\qty{0.01}{kg} が7こ\quad（2）\nobreak 407こ\quad（3）\nobreak 40.7、407、0.407}
\end{jibun}

\begin{reidai}[\R{小数}{しょう|すう}のある\R{量}{りょう}の\R{単位}{たん|い}をかえる]
\qty{1.25}{kg} を g で、\qty{350}{g} を kg で\R{表}{あらわ}しましょう。\par
\kangaekata $\qty{1}{kg}=\qty{1000}{g}$ です。kg を g にするには1000\R{倍}{ばい}、g を kg にするには1000でわります。$1000=10\times10\times10$ なので、\R{小数点}{しょう|すう|てん}は3つ\R{動}{うご}きます。\par
\shiki $\qty{1.25}{kg}=1.25\times\qty{1000}{g}$\par
\hspace*{2.6\zw}\mbox{1.25 → 12.5 → 125 → 1250}\quad（\R{右}{みぎ}へ3つ。\R{最後}{さい|ご}は0を\R{書}{か}く）\par
\hspace*{2.6\zw}\qty{1000}{g} で \qty{1}{kg} なので、350 を1000でわる\par
\hspace*{2.6\zw}\mbox{350 → 35 → 3.5 → 0.35}\quad（\R{左}{ひだり}へ3つ。\R{一}{いち}の\R{位}{くらい}に0を\R{書}{か}く）\par
\begin{bkcenter}
\begin{tikzpicture}[x=1mm, y=1mm]
  % 0〜2 kg と 0〜2000 g をならべた数直線（1 kg を 50 mm で）
  \foreach \yy/\u/\a/\b/\c in {6/kg/0/1/2, -6/g/0/1000/2000} {
    \draw[bkInk, line width=0.8pt] (0,\yy) -- (100,\yy);
    \foreach \k in {0,...,20} {\draw[bkInk, line width=0.4pt] (\k*5,\yy-1) -- (\k*5,\yy+1);}
    \node[anchor=west, font=\sffamily\footnotesize, text=bkGray, inner sep=0pt] at (102,\yy) {\u};}
  \foreach \x/\l in {0/0, 50/1, 100/2} {\node[anchor=south, font=\footnotesize, inner sep=0pt] at (\x,8.5) {\l};}
  \foreach \x/\l in {0/0, 50/1000, 100/2000} {\node[anchor=north, font=\footnotesize, inner sep=0pt] at (\x,-8.5) {\l};}
  \draw[bkOrange, line width=1.2pt] (62.5,7.5) -- (62.5,-7.5);
  \node[anchor=south, font=\sffamily\footnotesize, text=bkOrange!80!black, inner sep=0pt] at (62.5,13) {1.25};
  \node[anchor=north, font=\sffamily\footnotesize, text=bkOrange!80!black, inner sep=0pt] at (62.5,-13) {1250};
  \draw[bkBlue, line width=1.2pt] (17.5,7.5) -- (17.5,-7.5);
  \node[anchor=south, font=\sffamily\footnotesize, text=bkBlue!80!black, inner sep=0pt] at (17.5,13) {0.35};
  \node[anchor=north, font=\sffamily\footnotesize, text=bkBlue!80!black, inner sep=0pt] at (17.5,-13) {350};
\end{tikzpicture}
\end{bkcenter}
\kotae \qty{1250}{g}、\qty{0.35}{kg}
\end{reidai}

\begin{jibun}
\qty{2.4}{L} を mL で、\qty{85}{cm} を m で\R{表}{あらわ}しましょう。
\bkjans{\qty{2400}{mL}、\qty{0.85}{m}}
\end{jibun}

\begin{reidai}[\R{単位}{たん|い}をそろえてくらべる]
（1）\qty{1.5}{m} のロープと \qty{120}{cm} のロープ。\R{長}{なが}いのはどちらでしょう。\\
（2）\qty{0.9}{m}、\qty{95}{cm}、\qty{880}{mm} の3\R{本}{ぼん}のひもを、\R{長}{なが}い\R{順}{じゅん}にならべましょう。\par
\kangaekata \R{単位}{たん|い}を cm にそろえてから、\R{数値}{すう|ち}をくらべます。\R{数値}{すう|ち}だけをくらべてはいけません（150と120、1.5と120はくらべられない）。\par
\begin{bkcenter}
\begin{tikzpicture}[x=1mm, y=1mm]
  \bktape{(0,8)}{105}{bkBlue}{\qty{1.5}{m}（\qty{150}{cm}）}
  \bktape{(0,0)}{84}{bkOrange}{\qty{120}{cm}}
\end{tikzpicture}
\end{bkcenter}
\shiki （1）$\qty{1.5}{m}=1.5\times\qty{100}{cm}=\qty{150}{cm}$\par
\hspace*{4.6\zw}\qty{150}{cm} と \qty{120}{cm} では、150 のほうが\R{大}{おお}きい。\par
\hspace*{2.6\zw}（2）$\qty{0.9}{m}=0.9\times\qty{100}{cm}=\qty{90}{cm}$\par
\hspace*{4.6\zw}\qty{10}{mm} で \qty{1}{cm} なので、$880\div10=88$ で \qty{88}{cm}\par
\hspace*{4.6\zw}95、90、88 の\R{順}{じゅん}に\R{大}{おお}きい。\par
\kotae （1）\qty{1.5}{m} のロープ\quad（2）\qty{95}{cm}、\qty{0.9}{m}、\qty{880}{mm}
\end{reidai}

\begin{jibun}
（1）\qty{0.8}{kg} のメロンと \qty{750}{g} のメロン。\R{重}{おも}いのはどちらでしょう。\\
（2）\qty{2.3}{m}、\qty{225}{cm}、\qty{2400}{mm} を、\R{長}{なが}い\R{順}{じゅん}にならべましょう。
\bkjans{（1）\nobreak \qty{0.8}{kg}（\qty{800}{g}）のメロン\quad（2）\nobreak \qty{2400}{mm}、\qty{2.3}{m}、\qty{225}{cm}}
\end{jibun}

%% ---------------- まちがえやすいところ・物理工学 ----------------
\bksection{tsumazuki}{まちがえやすいところをチェック}

\begin{tsumazuki}{\R{答}{こた}えに\R{単位}{たん|い}をつけない}
\bkng{\qty{3}{m} は\R{何}{なん}cm？\quad こたえ\enspace 300}
\bkok{\qty{3}{m} は\R{何}{なん}cm？\quad こたえ\enspace \qty{300}{cm}}
「300」だけでは、cm なのか m なのかわかりません（\bkhatenaref の「あと5」と\R{同}{おな}じ）。
\end{tsumazuki}

\begin{tsumazuki}{かけるのか、わるのかを\R{逆}{ぎゃく}にする}
\bkng{$\qty{3}{m}=3\div\qty{100}{cm}=\qty{0.03}{cm}$}
\bkok{$\qty{3}{m}=3\times\qty{100}{cm}=\qty{300}{cm}$}
cm は m より\R{小}{ちい}さい\R{単位}{たん|い}なので、\R{数値}{すう|ち}は\R{大}{おお}きくなるはずです。\R{答}{こた}えの\R{大}{おお}きさを\R{確}{たし}かめよう。
\end{tsumazuki}

\begin{tsumazuki}{\R{単位}{たん|い}がちがうまま、\R{数値}{すう|ち}だけくらべる}
\bkng{\qty{1.5}{m} と \qty{120}{cm}。120 のほうが\R{大}{おお}きいので、\qty{120}{cm} が\R{長}{なが}い。}
\bkok{$\qty{1.5}{m}=\qty{150}{cm}$。150 と 120 で、\qty{1.5}{m} が\R{長}{なが}い。}
1.5 は「\qty{1}{m} の1.5こ\R{分}{ぶん}」、120 は「\qty{1}{cm} の120こ\R{分}{ぶん}」。もとにする\R{大}{おお}きさ（\R{単位}{たん|い}）がちがう\R{数値}{すう|ち}は、そのままではくらべられません。
\end{tsumazuki}

\begin{tsumazuki}{\R{小数点}{しょう|すう|てん}を\R{動}{うご}かすとき、0をわすれる}
\bkng{$\qty{0.4}{m}=\qty{4}{cm}$\quad（\R{右}{みぎ}へ1つしか\R{動}{うご}かしていない）}
\bkok{$\qty{0.4}{m}=0.4\times\qty{100}{cm}=\qty{40}{cm}$\quad（\R{右}{みぎ}へ2つ）}
100\R{倍}{ばい}は\R{右}{みぎ}へ2つ。\R{数字}{すう|じ}のない\R{位}{くらい}には0を\R{書}{か}きます（\mbox{0.4 → 4 → 40}）。
\end{tsumazuki}

\begin{butsuri}
\noindent\begin{minipage}[c]{0.42\linewidth}\centering
\begin{tikzpicture}[x=1mm, y=1mm]
  \bkiibshelf{1.8}{800}{1800}{6}
\end{tikzpicture}
\end{minipage}\hfill
\begin{minipage}[c]{0.56\linewidth}
\textbf{\R{設計図}{せっ|けい|ず}の\R{寸法}{すん|ぽう}}\par
\R{家具}{か|ぐ}や\R{機械}{き|かい}の\R{設計図}{せっ|けい|ず}では、\R{長}{なが}さをすべて mm で\R{書}{か}き、\R{単位}{たん|い}は\R{書}{か}かないやくそくです。\R{図}{ず}の1800は、$\qty{1800}{mm}=\qty{1.8}{m}$。\R{単位}{たん|い}をそろえておくので、だれが\R{読}{よ}んでも、まちがえずに\R{部品}{ぶ|ひん}を\R{作}{つく}れます。
\end{minipage}

\par\vspace{1.5mm}
\noindent\begin{minipage}[c]{0.36\linewidth}\centering
\begin{tikzpicture}[x=1mm, y=1mm]
  % デジタルはかり
  \fill[bkGray!35, rounded corners=2mm] (0,0) rectangle (42,20);
  \fill[bkGray!55, rounded corners=1mm] (4,20) rectangle (38,23);
  \bkdisplay{(5,4.5)}{32}{0.350\,kg}
\end{tikzpicture}
\end{minipage}\hfill
\begin{minipage}[c]{0.61\linewidth}
\textbf{\R{計器}{けい|き}の\R{表示}{ひょう|じ}}\par
\R{電子}{でん|し}はかりの「0.350\,kg」と「350\,g」は\R{同}{おな}じ\R{量}{りょう}です。\R{計器}{けい|き}を\R{読}{よ}むときは、まず\R{表示}{ひょう|じ}の\R{単位}{たん|い}を\R{見}{み}ます。\R{単位}{たん|い}を\R{見}{み}\R{落}{お}とすと、1000\R{倍}{ばい}もまちがえてしまいます。
\end{minipage}
\end{butsuri}

\begin{column}{\R{昔}{むかし}の\R{日本}{に|ほん}の\R{長}{なが}さの\R{単位}{たん|い}}
「\R{一寸法師}{いっ|すん|ぼう|し}」の1\R{寸}{すん}は\R{約}{やく} \qty{3}{cm}、10\R{寸}{すん}の1\R{尺}{しゃく}は\R{約}{やく} \qty{30}{cm}。\R{竹}{たけ}の\R{笛}{ふえ}「\R{尺八}{しゃく|はち}」は、1\R{尺}{しゃく}8\R{寸}{すん}（\R{約}{やく} \qty{55}{cm}）の\R{長}{なが}さから\R{名前}{な|まえ}がつきました。\R{単位}{たん|い}がかわっても、\R{長}{なが}さそのものは\R{同}{おな}じです。
\end{column}


\clearpage
%% ---------------- 章末問題 ----------------
\bkendproblems
\bkproblemgroup{\R{基本}{き|ほん}}{1}

\bkproblem \R{次}{つぎ}の\R{量}{りょう}は、かぞえる\R{数}{かず}・はかる\R{数}{かず}のどちらで\R{表}{あらわ}すでしょう。\\
\mbox{（1）\R{体重}{たい|じゅう}}\quad \mbox{（2）\R{本}{ほん}のページ\R{数}{すう}}\quad \mbox{（3）\R{教室}{きょう|しつ}の\R{温度}{おん|ど}}\quad \mbox{（4）\R{自転車}{じ|てん|しゃ}の\R{台数}{だい|すう}}
\bkpans{（1）\nobreak はかる\R{数}{かず}\quad（2）\nobreak かぞえる\R{数}{かず}\quad（3）\nobreak はかる\R{数}{かず}\quad（4）\nobreak かぞえる\R{数}{かず}}

\bkproblem \R{次}{つぎ}の\R{量}{りょう}を、［\quad］の\R{単位}{たん|い}で\R{表}{あらわ}しましょう。\\
（1）\qty{6}{m}［cm］\quad（2）\qty{5}{cm}［mm］\quad（3）\qty{3000}{mL}［L］\quad（4）\qty{2}{kg}［g］
\bkpans{（1）\nobreak \qty{600}{cm}\quad（2）\nobreak \qty{50}{mm}\quad（3）\nobreak \qty{3}{L}\quad（4）\nobreak \qty{2000}{g}}

\bkproblem \qty{3.82}{m} の 8 と 2 は、それぞれ\R{何}{なに}が\R{何}{なん}こあることを\R{表}{あらわ}していますか。また、3.82 は 0.01 を\R{何}{なん}こ\R{集}{あつ}めた\R{数}{かず}ですか。
\bkpans{8 は \qty{0.1}{m} が8こ、2 は \qty{0.01}{m} が2こ。382こ}

\bkproblem 0.57 と 4.3 を、それぞれ10\R{倍}{ばい}、100\R{倍}{ばい}した\R{数}{かず}と、$\dfrac{1}{10}$にした\R{数}{かず}を\R{求}{もと}めましょう。
\bkpans{0.57 → 5.7、57、0.057\quad \mbox{4.3 → 43}、430、0.43}

\bkproblem \R{次}{つぎ}の\R{量}{りょう}を、［\quad］の\R{単位}{たん|い}で\R{表}{あらわ}しましょう。\\
（1）\qty{2.6}{m}［cm］\quad（2）\qty{0.75}{L}［mL］\quad（3）\qty{480}{g}［kg］\quad（4）\qty{35}{mm}［cm］
\bkpans{（1）\nobreak \qty{260}{cm}\quad（2）\nobreak \qty{750}{mL}\quad（3）\nobreak \qty{0.48}{kg}\quad（4）\nobreak \qty{3.5}{cm}}

\bkproblem \R{図}{ず}のえんぴつの\R{長}{なが}さを、（1）○cm○mm\quad（2）cm だけ\quad（3）mm だけ で\R{表}{あらわ}しましょう。
\begin{bkcenter}
\begin{tikzpicture}[x=1mm, y=1mm]
  \bkiibruler{(0,0)}{0.8}{15}
  \bkpencil{(2.2,3.6)}{98.6}{bkGreen}
\end{tikzpicture}
\end{bkcenter}
\bkpans{（1）\nobreak 12\,cm\,6\,mm\quad（2）\nobreak \qty{12.6}{cm}\quad（3）\nobreak \qty{126}{mm}}

\clearpage
\bkproblemgroup{チャレンジ}{1}

\bkproblem \qty{2}{m} のリボンを、\qty{1}{cm} ずつに\R{切}{き}ると、\R{何本}{なん|ぼん}になりますか。\qty{1}{mm} ずつに\R{切}{き}ると、\R{何本}{なん|ぼん}になりますか。
\bkpans{200\R{本}{ぽん}、2000\R{本}{ぼん}}

\bkproblem （1）\qty{2.1}{m} と \qty{205}{cm} では、どちらが\R{長}{なが}いでしょう。（2）\qty{1.3}{m}、\qty{128}{cm}、\qty{1350}{mm} の3\R{本}{ぼん}のひもを、\R{長}{なが}い\R{順}{じゅん}にならべましょう。
\bkpans{（1）\nobreak \qty{2.1}{m}（\qty{210}{cm}）\quad（2）\nobreak \qty{1350}{mm}、\qty{1.3}{m}、\qty{128}{cm}（cm にそろえると 135、130、128）}

\noindent\begin{minipage}[c]{0.62\linewidth}
\bkproblem \R{右}{みぎ}の\R{本}{ほん}だなの\R{設計図}{せっ|けい|ず}（\R{単位}{たん|い}は mm）を\R{見}{み}て\R{答}{こた}えましょう。\\
（1）\R{高}{たか}さは\R{何}{なん}m ですか。\\（2）たなとたなのあいだ（250）は\R{何}{なん}cm ですか。\\（3）\R{横}{よこ}はばは\R{何}{なん}m ですか。
\end{minipage}\hfill
\begin{minipage}[c]{0.36\linewidth}\centering
\begin{tikzpicture}[x=1mm, y=1mm]
  \bkiibshelf{2.6}{600}{1500}{6}
\end{tikzpicture}
\end{minipage}
\bkpans{（1）\nobreak \qty{1.5}{m}\quad（2）\nobreak \qty{25}{cm}\quad（3）\nobreak \qty{0.6}{m}}

\begin{furikaeri}
\begin{checklist}
\item 「かぞえる\R{数}{かず}」と「はかる\R{数}{かず}」のちがいを\R{説明}{せつ|めい}できる。\hfill{\footnotesize\sffamily\color{bkTheme}→\pageref{bk:2-1:kazoeru}ページ}
\item \R{量}{りょう}を「\R{数値}{すう|ち}×\R{単位}{たん|い}」で\R{表}{あらわ}し、\R{単位}{たん|い}をかえて\R{表}{あらわ}せる。\hfill{\footnotesize\sffamily\color{bkTheme}→\pageref{bk:2-1:tani}ページ}
\item \R{小数}{しょう|すう}の\R{位}{くらい}がわかり、10\R{倍}{ばい}・$\dfrac{1}{10}$ で\R{小数点}{しょう|すう|てん}を\R{動}{うご}かせる。\hfill{\footnotesize\sffamily\color{bkTheme}→\pageref{bk:2-1:kurai}ページ}
\end{checklist}
チェックがつかなかったところは、→のページにもどって、もう\R{一度}{いち|ど}\R{読}{よ}んでみよう。
\end{furikaeri}

\bkfukidashi{mina}{\R{答}{こた}えに\R{単位}{たん|い}をつけると、\R{大}{おお}きさが\R{合}{あ}っているかも\R{確}{たし}かめられるね。}
